% ABET Student Outcome Mapping
% Auto-generated from megn300_master_reference.tex

\chapter*{ABET Student Outcome Mapping}
\addcontentsline{toc}{chapter}{ABET Student Outcome Mapping}

MEGN~300 is designed to develop the following ABET student outcomes. The table below maps each outcome to the course activities and document chapters that most directly address it.

\begin{center}
\begin{longtable}{p{0.8cm}>{\raggedright\arraybackslash}p{5.5cm}>{\raggedright\arraybackslash}p{5.5cm}}
\toprule
\textbf{SO} & \textbf{ABET Student Outcome} & \textbf{Course Mapping} \\
\midrule
\endfirsthead
\toprule
\textbf{SO} & \textbf{ABET Student Outcome} & \textbf{Course Mapping} \\
\midrule
\endhead
1 & Identify, formulate, and solve complex engineering problems by applying principles of engineering, science, and mathematics. & Measurement theory (Ch~1--3), error analysis and uncertainty propagation (Ch~4), signal processing (Ch~6--12), PID control design (Ch~15--17), heat transfer modeling (Ch~18). Labs require formulating measurement problems and applying mathematical models to physical systems. \\
\midrule
2 & Apply engineering design to produce solutions that meet specified needs with consideration of public health, safety, and welfare, as well as global, cultural, social, environmental, and economic factors. & Safety chapter (Laboratory Safety), circuit protection design (Ch~14--15), functional safety requirements (Ch~15), sensor selection for real-world constraints, energy-efficient signal conditioning design. \\
\midrule
3 & Communicate effectively with a range of audiences. & Lab reports requiring clear documentation of methodology, results, and uncertainty. Practice questions requiring written technical explanations. Requirements mapping boxes throughout demonstrate professional specification language. \\
\midrule
5 & Function effectively on a team whose members together provide leadership, create a collaborative and inclusive environment, establish goals, plan tasks, and meet objectives. & Lab exercises conducted in teams. Integration of multiple subsystems (sensors, DAQ, signal conditioning, control) requiring coordinated team effort and interface management. \\
\midrule
6 & Develop and conduct appropriate experimentation, analyze and interpret data, and use engineering judgment to draw conclusions. & All lab exercises involve experimental design, data collection, statistical analysis (Ch~5), FFT-based frequency analysis (Ch~8--9), uncertainty quantification (Ch~4), and engineering judgment in interpreting results. This is the core competency of the course. \\
\midrule
7 & Acquire and apply new knowledge as needed, using appropriate learning strategies. & Students must independently read datasheets, select sensors, configure DAQ hardware, learn LabVIEW programming (Appendix), and apply filter design tools. The Master Reference Document serves as a self-study resource supporting independent learning. \\
\bottomrule
\end{longtable}
\end{center}

\textbf{Note}: ABET Student Outcome~4 (recognizing ethical and professional responsibilities) is addressed primarily through the safety practices, proper uncertainty reporting (honest representation of measurement quality), and the emphasis on documentation standards throughout the course. These outcomes are assessed through lab reports, exams, and the final project deliverables.

